% GENERATED FILE. Edit canonical lessons, then run npm run generate:books.
% canonical-source-hash: fnv1a64:3596eccb70d8e7fc
% canonical-lessons: FA-C14-khahar, FA-C14-pesar, FA-C14-mard, FA-C14-zan, FA-C14-dust

\chapter{Sister, Son, Man, Woman, Friend}
\label{ch:fa-people-words}
\hlchaptermodality{14}

\begin{chapteropening}
\textbf{By the end of this chapter:} \emph{I can name my sister, son, and friend in Persian, tell man and woman apart, and open and close a short interaction with all of them.}

The chapter and the tranche close on dust, friend, then run all five people-words back to back and rebuild a complete short interaction with a friend: a polite request for tea, a request for a key, and the standard farewell both voices share.
\end{chapteropening}

\section[\texorpdfstring{\fa{خواهر} (khâhar)}{خواهر (khâhar)}]{\fa{خواهر} — sister, the family word the earlier tranche skipped}

\label{lesson:FA-C14-khahar}



\begin{quote}
Say \emph{bârân}, and whether it has an English cousin. The sky words are done; the family the earlier tranche built — \emph{mâdar, pedar, barâdar, dokhtar} — is missing exactly one member.
\end{quote}

\subsection*{You'll want to know first — one word, a familiar silent letter}
\begin{quote}
\textbf{\fa{خواهر}} — \emph{khâhar} — \textbf{sister}
\end{quote}

Five letters, all already yours: \textbf{\fa{خ} \fa{و} \fa{ا} \fa{ه} \fa{ر}}. The \textbf{\fa{و}} here is silent, exactly the pattern \emph{khândan} taught: written after \textbf{\fa{خ}} and before \textbf{\fa{ا}}, never pronounced. Say it \textbf{kh-â-har}.

\begin{cousinweb}
\textbf{\fa{خواهر}} continues Middle Persian \textbf{xwāhar}, from Indo-European *\emph{swésōr}, \textquotedblleft{}sister\textquotedblright{} — as secure a cognate as \emph{barâdar}'s own *\emph{b\textsuperscript{h}réh\textsubscript{2}tēr}. English kept it as \textbf{sister} (through Old Norse, replacing the native Old English form), Latin as \textbf{soror} ($\to$ \textbf{sorority}), and Sanskrit as \textbf{svasṛ}. \emph{mâdar, pedar, barâdar, dokhtar, khâhar} now give the track all five of Indo-European's plainest kinship words.
\end{cousinweb}

\subsection*{Your turn}
\begin{itemize}
\raggedright
  \item \emph{Say it:} \textbf{khâhar} — sister
  \item \emph{Connect:} \textbf{khâhar} $\leftarrow$ *\emph{swésōr} $\to$ English \textbf{sister}, Latin \textbf{soror}, Sanskrit \textbf{svasṛ}
  \item \emph{Read it:} the silent \textbf{\fa{و}} in \textbf{\fa{خواهر}}, the same pattern \emph{khândan} taught
  \item \emph{Name:} all five family words this track now holds — \emph{mâdar, pedar, barâdar, dokhtar, khâhar}
  \item \emph{Retrieve:} \emph{ketâb} — book, from a few lessons back
\end{itemize}

\begin{tcolorbox}[breakable,colback=teal!4,colframe=teal!35!black,title={Before you move on}]
What does \textbf{\fa{خواهر}} mean? (\textbf{Sister.}) Which letter in it is silent, and which earlier word taught the same pattern? (\textbf{\fa{و}}; \emph{khândan}.) Which English word is its direct cousin? (\textbf{Sister.})

Source: \href{https://en.wiktionary.org/wiki/\%D8\%AE\%D9\%88\%D8\%A7\%D9\%87\%D8\%B1}{Wiktionary: \fa{خواهر}}.
\end{tcolorbox}
\section[\texorpdfstring{\fa{پسر} (pesar)}{پسر (pesar)}]{\fa{پسر} — son, \emph{dokhtar}'s missing counterpart}

\label{lesson:FA-C14-pesar}



\begin{quote}
Say \emph{khâhar}, and its secure English cousin. Then say \emph{barâdar} — its own family tranche taught \emph{dokhtar}, \textquotedblleft{}daughter,\textquotedblright{} but never gave her brother's counterpart a lesson of its own, until now.
\end{quote}

\subsection*{You'll want to know first — one word}
\begin{quote}
\textbf{\fa{پسر}} — \emph{pesar} — \textbf{son}
\end{quote}

Three letters, all already yours: \textbf{\fa{پ} \fa{س} \fa{ر}} — \textbf{\fa{پ}} \emph{pe}, from \emph{porsidan}, then \emph{s}, \emph{r}.

\begin{cousinweb}
\textbf{\fa{پسر}} continues Old Persian \textbf{puça}, cognate with Sanskrit \textbf{putra}, \textquotedblleft{}son\textquotedblright{} — an Indo-Iranian word Persian and Sanskrit still share with each other. English \textbf{son} is not part of this family: it descends from an entirely different Indo-European root, so — as with \emph{bârân} and \emph{raftan} before it — no English cousin is claimed here. \textbf{\fa{پسر}} and \emph{dokhtar} now give the track a complete pair for \textquotedblleft{}son\textquotedblright{} and \textquotedblleft{}daughter,\textquotedblright{} even though only one of the two has an English relative.
\end{cousinweb}

\subsection*{Your turn}
\begin{itemize}
\raggedright
  \item \emph{Say it:} \textbf{pesar} — son
  \item \emph{Connect:} \textbf{pesar} $\leftarrow$ Old Persian \textbf{puça} $\to$ Sanskrit \textbf{putra} — no English cousin
  \item \emph{Run through:} \textbf{barâdar, pesar} — brother, son — the two words this pair now completes
  \item \emph{Run through:} \textbf{khâhar, pesar} — sister, son
\end{itemize}

\begin{tcolorbox}[breakable,colback=teal!4,colframe=teal!35!black,title={Before you move on}]
What does \textbf{\fa{پسر}} mean? (\textbf{Son.}) Which Sanskrit word is its direct cousin? (\textbf{Putra}.) Does English \textbf{son} share the same root? (\textbf{No.})

Source: \href{https://en.wiktionary.org/wiki/\%D9\%BE\%D8\%B3\%D8\%B1}{Wiktionary: \fa{پسر}}.
\end{tcolorbox}
\section[\texorpdfstring{\fa{مرد} (mard)}{مرد (mard)}]{\fa{مرد} — man, but its root says \textquotedblleft{}mortal\textquotedblright{}}

\label{lesson:FA-C14-mard}



\begin{quote}
Say \emph{pesar}, and the Sanskrit word it shares a root with. This lesson's word looks like it should be an easy English cousin. It is not.
\end{quote}

\subsection*{You'll want to know first — one word}
\begin{quote}
\textbf{\fa{مرد}} — \emph{mard} — \textbf{man}
\end{quote}

Three letters, all already yours: \textbf{\fa{م} \fa{ر} \fa{د}} — \emph{m}, \emph{r}, \emph{d}.

\begin{cousinweb}
\textbf{\fa{مرد}} continues Middle Persian \textbf{mard}, from Indo-European *\emph{mer-}, \textquotedblleft{}to die.\textquotedblright{} Its literal sense is not \textquotedblleft{}man\textquotedblright{} at all — it is \textquotedblleft{}mortal,\textquotedblright{} \textquotedblleft{}one who dies.\textquotedblright{} The same root gives Latin \textbf{mortuus}, \textquotedblleft{}dead\textquotedblright{} ($\to$ English \textbf{mortal}, \textbf{mortician}), Greek \textbf{brotós} ($\to$ \textbf{ambrosia}, \textquotedblleft{}not mortal\textquotedblright{}), and Sanskrit \textbf{martya}. English \textbf{man} looks like an obvious match for \textbf{\fa{مرد}} and is not one — it descends from a completely different root — so this is one place where guessing a cousin from sound alone would mislead rather than help.
\end{cousinweb}

\subsection*{Your turn}
\begin{itemize}
\raggedright
  \item \emph{Say it:} \textbf{mard} — man
  \item \emph{Connect:} \textbf{mard} $\leftarrow$ *\emph{mer-} $\to$ Latin \textbf{mortuus}, \textbf{mortal}; Greek \textbf{brotós}; Sanskrit \textbf{martya} — not English \textbf{man}
  \item \emph{Run through:} \textbf{pesar, mard} — son, man
\end{itemize}

\begin{tcolorbox}[breakable,colback=teal!4,colframe=teal!35!black,title={Before you move on}]
What does \textbf{\fa{مرد}} mean? (\textbf{Man.}) What does its root literally mean, and which English word keeps that sense? (\textbf{Mortal}; \textbf{mortal} itself.) Does English \textbf{man} share the same root? (\textbf{No — a false friend.})

Source: \href{https://en.wiktionary.org/wiki/\%D9\%85\%D8\%B1\%D8\%AF}{Wiktionary: \fa{مرد}}.
\end{tcolorbox}
\section[\texorpdfstring{\fa{زن} (zan)}{زن (zan)}]{\fa{زن} — woman, a root hiding behind a royal title}

\label{lesson:FA-C14-zan}



\begin{quote}
Say \emph{mard}, and why it is not a cousin of English \textbf{man}. This lesson's word finds its English cousin in an unexpected place.
\end{quote}

\subsection*{You'll want to know first — one word}
\begin{quote}
\textbf{\fa{زن}} — \emph{zan} — \textbf{woman}
\end{quote}

Two letters, both already yours: \textbf{\fa{ز} \fa{ن}} — \textbf{\fa{ز}} \emph{ze}, from \emph{zabân}, then \emph{n}.

\begin{cousinweb}
\textbf{\fa{زن}} continues Middle Persian \textbf{zan}, from Indo-European *\emph{g\textsuperscript{w}\'{\={e}}n(h\textsubscript{2})-}, \textquotedblleft{}woman.\textquotedblright{} English kept the same root, but only inside a narrowed, elevated word: \textbf{queen} was once simply \textquotedblleft{}woman,\textquotedblright{} the way \emph{mard} narrowed \textquotedblleft{}mortal\textquotedblright{} down to \textquotedblleft{}man.\textquotedblright{} Greek kept the plain sense directly as \textbf{gynē} ($\to$ \textbf{gynecology}), and Russian as \textbf{žena}. \emph{mard} and \textbf{\fa{زن}} give the track its first pair of person words that are not family roles.
\end{cousinweb}

\subsection*{Your turn}
\begin{itemize}
\raggedright
  \item \emph{Say it:} \textbf{zan} — woman
  \item \emph{Connect:} \textbf{zan} $\leftarrow$ *\emph{g\textsuperscript{w}\'{\={e}}n(h\textsubscript{2})-} $\to$ English \textbf{queen}, Greek \textbf{gynē}
  \item \emph{Run through:} \textbf{mard, zan} — man, woman
  \item \emph{Contrast:} \emph{mard} narrowed from \textquotedblleft{}mortal\textquotedblright{}; \textbf{\fa{زن}} widened English's cousin down to royalty alone
  \item \emph{Retrieve:} \emph{setâre} — star, one more well-attested cousin
\end{itemize}

\begin{tcolorbox}[breakable,colback=teal!4,colframe=teal!35!black,title={Before you move on}]
What does \textbf{\fa{زن}} mean? (\textbf{Woman.}) Which English word is its cousin, and what did that word originally mean before it narrowed? (\textbf{Queen}; simply \textquotedblleft{}woman.\textquotedblright{}) Which Greek word shares the same root? (\textbf{Gynē}.)

Source: \href{https://en.wiktionary.org/wiki/\%D8\%B2\%D9\%86}{Wiktionary: \fa{زن}}.
\end{tcolorbox}
\section[\texorpdfstring{\fa{دوست} (dust)}{دوست (dust)}]{\fa{دوست} — friend, and everything this tranche has built}

\label{lesson:FA-C14-dust}



\begin{quote}
Say \emph{mard, zan} — man, woman. One more person word closes the tranche, and it is a word this track has half-taught already.
\end{quote}

\subsection*{You'll want to know first — one word, already half-known}
\begin{quote}
\textbf{\fa{دوست}} — \emph{dust} — \textbf{friend}
\end{quote}

Four letters, all already yours: \textbf{\fa{د} \fa{و} \fa{س} \fa{ت}} — \emph{d}, long \emph{u}, \emph{s}, \emph{t}. The compound verb for \textquotedblleft{}to love\textquotedblright{} already ran on this exact word; here it stands alone, as the person it names.

\begin{cousinweb}
\textbf{\fa{دوست}} goes back through Old Persian \textbf{dauštā} to Indo-European *\emph{ǵews-}, \textquotedblleft{}to taste, to choose\textquotedblright{} — a friend is \emph{the one chosen}. English inherited the root directly as \textbf{choose}; Latin took it as \textbf{gustus}, \textquotedblleft{}taste,\textquotedblright{} giving English \textbf{gusto} and \textbf{disgust}. The compound-verb lesson that first taught this root only ever used \textbf{\fa{دوست}} glued to \emph{dâshtan}; this is its first lesson as a free-standing noun.
\end{cousinweb}

\subsection*{Your turn — the tranche, start to finish}
One short scene, built entirely from independently learned pieces:

\begin{itemize}
\raggedright
  \item \emph{Say it:} \textbf{dust} — friend; then \textbf{khâhar, pesar, mard, zan, dust} — all five people-words this closing chapter taught
  \item \emph{Connect:} \textbf{dust} $\leftarrow$ *\emph{ǵews-} $\to$ English \textbf{choose}, Latin \textbf{gustus}
  \item \emph{Say it:} \emph{chây, lotfan} — offer your friend tea in the \textbf{bârân}, rain, reusing the very first polite request this track built
  \item \emph{Say it:} \emph{kelid-e man, lotfan} — and ask a friend to pass your key
  \item \emph{Say it:} \emph{hâfez} — guardian, protector — the Arabic-rooted half of a farewell you have used since your very first complete exchange
  \item \emph{Say it:} joined \emph{khodâ hâfez}, and recall what it leaves unsaid — no verb, just two halves fused into one wish
  \item \emph{Run through:} \emph{Salâm! Hâl-e shomâ chetor ast? … Khubam, mamnun. … Khodâ hâfez. … Khodâ hâfez.} — the exact four-line interaction a friend opens and closes
\end{itemize}

\begin{tcolorbox}[breakable,colback=teal!4,colframe=teal!35!black,title={Before you move on}]
What does \textbf{\fa{دوست}} mean on its own? (\textbf{Friend.}) Which compound verb already ran on this exact word? (\textbf{Dust dâshtan}, \textquotedblleft{}to love.\textquotedblright{}) Name all five people-words this closing chapter taught, in order. (\textbf{Khâhar, pesar, mard, zan, dust.}) How does a friend close the four-line interaction? (\emph{khodâ hâfez}, said by both voices.)

Source: \href{https://en.wiktionary.org/wiki/\%D8\%AF\%D9\%88\%D8\%B3\%D8\%AA}{Wiktionary: \fa{دوست}}.
\end{tcolorbox}
